\section{Conclusão}
\label{sec:conclusao}

Apresentamos a segmentação dos padrões de Turing do Gray-Scott segundo o protocolo comum da turma.
O GT é físico ($B>\tau$, congelado antes da comparação), e a entrada é restrita a $A$ para evitar
vazamento de rótulo. Usamos todas as 1200 trajetórias da divisão oficial, resolução nativa
$128\times128$ sem ampliação, 20 épocas com checkpoint da melhor época de validação, operações
periódicas e seleção de parâmetros só na validação. A análise mecanística prevê quatro resultados:
\begin{enumerate}
  \item normalizações afins (P1--P4) são neutras na U-Net, por causa da normalização em lote, mas
    podem afetar as redes pré-treinadas, que esperam entradas na escala do ImageNet;
  \item filtros espaciais (P6--P10) não ajudam em dados de simulação sem ruído, e o gaussiano piora
    a fronteira;
  \item a U-Net, que prediz na resolução completa, supera SegFormer e Mask2Former nas estruturas
    pequenas, cujos erros se concentram nas bordas reconstruídas a partir de $1/4$ da resolução;
  \item a filtragem temporal (Q5) só é segura nos regimes estáticos.
\end{enumerate}
No piloto, (2), (3) e (4) se confirmaram: nenhum filtro melhorou de forma significativa, a U-Net
superou os Transformers em 15 a 17 pontos de Dice, com erros deles concentrados em pontes de FP
entre estruturas vizinhas, e Q5 só foi inócuo em Maze e Bubbles. A hipótese (1) foi refinada: o
Z-score mudou a \emph{calibração}, e não a separabilidade. Com o limiar otimizado (Q1), ele produziu
a melhor configuração da U-Net (Dice 0,972). Por isso, no benchmark, o efeito de cada
pré-processamento deve ser lido junto com Q1.
\revisar{Atualizar com os números finais do experimento completo.} As tabelas estão no esquema
padronizado do protocolo, prontas para o artigo de benchmark entre sistemas físicos.
